% ============================================================
% COMPOSER'S NOTE § 3 — A Trombone in Third Position
% ============================================================
% The origin story: the band director, the Yamaha vs. the Bach,
% the specific sharpness of third position on one instrument
% vs. another. The epistemological lesson:
% instruments carry the intentions and sonic ideals of their makers.
% This is the first sounding of the leitmotif — the drone
% surfaces into consciousness here, then goes back under.
%
% Anchor sentence: "To modify a tool is to engage in an epistemic act."
%
% No external citations needed here — this is the researcher's
% own story. The theory is implicit.
% ============================================================

\section*{A Trombone in Third Position}

In high school, as a young musician, it was explained to me by my assistant band director that my trombone has a specific temperament. My Yamaha produced a slightly sharper tone in the typical ``third position'' area of the slide (when tuned to a standard pitch), different from that of a Bach or Shire model. Each maker, he said, designed, built, and tuned the instrument according to a distinct philosophy of sound. That detail---which is just mere millimeters of brass---taught me that instruments carry with them the intentions and sonic ideals of their makers. Now, I recognize that moment as a lesson in epistemology: every device, from a trombone to a software synthesizer, expresses a particular worldview about what counts as correct, beautiful, or possible.

This dissertation begins from such recognition. To modify a tool is to engage in an epistemic act: to re-tune inherited assumptions and redirect the flow of knowledge. Through the framework of \textit{Speculocultural Technopoiesis}, I treat modification as a process of transduction---where speculative cultural imagination takes material form---and refraction---where altered instruments bend perception toward new sonic and cultural relationships. The research situates these acts within the intertwined legacies of \textit{sonic virtuality} \parencite{Grimshaw_Garner_2015}, \textit{sonic fiction} \parencite{Eshun_1999}, and culturally grounded practices of making, positioning the sonic as both medium and method for re-imagining technological knowledge.

\adnote{Extracted verbatim from \texttt{intro\_context.tex}---this is your prose and needs the least revision of any section. The citations were already present. Consider whether you want to expand the trombone memory itself before the epistemological turn, or whether this length and pacing is right.}
